% Part: sets-functions-relations
% Chapter: sets
% Section: russells-paradox

\documentclass[../../../include/open-logic-section]{subfiles}


\begin{document}

\olfileid{sfr}{set}{rus}
\olsection{De paradox van Russell}

Dankzij extensionaliteit mogen we $\Setabs{x}{\phi(x)}$ schrijven voor
\emph{de} verzameling van alle $x$ waarvoor~$\phi(x)$ geldt. Maar uit
extensionaliteit volgt \emph{eigenlijk} alleen dit. \emph{Als} er een
verzameling bestaat met precies de objecten waarvoor~$\phi$ geldt als
elementen, \emph{dan} bestaat er maar één zo'n verzameling. Anders
gezegd: zodra een bepaalde~$\phi$ vaststaat, is
$\Setabs{x}{\phi(x)}$ uniek, \emph{als die verzameling bestaat}.

Maar dat ‘als die verzameling bestaat’ is belangrijk!{} Bij
\emph{comprehensie} vorm je een verzameling uit alle objecten met een
bepaalde eigenschap. Dat kan niet bij elke eigenschap. Sommige
eigenschappen bepalen dus \emph{geen} verzameling. Als dat wel kon,
zouden we direct in tegenspraak komen. Het bekendste voorbeeld is de
paradox van Russell.

Verzamelingen kunnen zelf !!{element}s van andere verzamelingen
zijn---de machtsverzameling van een verzameling~$A$ bestaat
bijvoorbeeld uit verzamelingen. Daarom kunnen we vragen of een
verzameling !!a{element} van een andere verzameling is. Kan een
verzameling element van zichzelf zijn? Niets aan het idee van een
verzameling lijkt dit onmogelijk te maken. Stel bijvoorbeeld dat
\emph{alle} verzamelingen samen één groep objecten vormen. Dan zou je
kunnen denken dat je daar één verzameling van kunt maken---de
verzameling van alle verzamelingen. Omdat die zelf ook een verzameling
is, zou zij !!a{element} zijn van de verzameling van alle
verzamelingen.

De paradox van Russell begint bij de eigenschap dat een verzameling
zichzelf niet als !!a{element} heeft, dus \emph{geen element van
zichzelf is}. Stel nu dat er een verzameling bestaat van alle
verzamelingen die zichzelf niet als !!a{element} hebben. Bestaat
\[
R = \Setabs{x}{x \notin x}
\]
dan? We kunnen bewijzen dat deze verzameling niet bestaat.

\begin{thm}[De paradox van Russell]\ollabel{thm:russells-paradox}
	Er bestaat geen verzameling $R = \Setabs{x}{x \notin x}$.
\end{thm}

\begin{proof}
Als $R = \Setabs{x}{x \notin x}$ bestaat, dan geldt
$R \in R$ precies wanneer $R \notin R$. Dat spreekt zichzelf tegen.
\end{proof}

\begin{tagblock}{novice}
\begin{explain}
We nemen dit bewijs stap voor stap door. Als $R$ bestaat, kunnen we
vragen of $R \in R$ geldt. Neem eerst aan dat $R \in R$. De
verzameling~$R$ bevat volgens haar definitie alle verzamelingen die
geen !!{element}s van zichzelf zijn. Maar als $R \in R$, is $R$ wel
element van zichzelf. Dan mist die verzameling dus de eigenschap waarmee $R$
is gedefinieerd. Alleen verzamelingen met die eigenschap horen
bij~$R$. Daarom kan $R$ geen !!{element} van~$R$ zijn: $R \notin R$.
Maar $R$ kan niet tegelijk wel en niet !!a{element} van~$R$ zijn.
Dat is een tegenspraak.

De aanname $R \in R$ geeft een tegenspraak, dus geldt $R \notin R$.
Maar ook dat geeft een tegenspraak!{} Als $R \notin R$, heeft $R$
juist wel de eigenschap waarmee $R$ is gedefinieerd. Dan zou $R$, net
als alle andere verzamelingen die geen element van zichzelf zijn, wel
!!a{element} van~$R$ zijn. Opnieuw kan $R$ niet tegelijk geen en wel
!!a{element} van zichzelf zijn.
\end{explain}
\end{tagblock}

\begin{digress}
Hoe bouwen we een verzamelingenleer waarin we niet in de paradox van
Russell terechtkomen? We mogen daarin dus niet de
\emph{tegenstrijdige} bewering doen dat $R = \Setabs{x}{x \notin x}$
bestaat. We hebben daarom axioma's nodig: regels die heel precies
zeggen wanneer verzamelingen bestaan (en wanneer niet).
	
De verzamelingenleer in dit hoofdstuk doet dat niet. Ze is
\emph{echt na\"ief}. Ze zegt alleen dat verzamelingen aan
extensionaliteit voldoen en dat je, als je al enkele verzamelingen
hebt, hun vereniging, doorsnede enzovoort kunt maken. We kunnen de
verzamelingenleer ook nauwkeuriger opbouwen, met precieze regels.
\oliflabeldef{cumul:::part}{Die strengere versie komt aan bod in Deel
\olref[cumul][][]{part}. Voorlopig werken we na\"ief verder en letten
we goed op dat we geen tegenspraken krijgen.}{}
\end{digress}
\end{document}
